\chapter{Negócios, Gestão e Comunicação Profissional}
\label{ch:negocios}

\ficcao Na nota da reunião lê-se «ver se há sala na quinta-feira». No resumo aparece «reunião marcada para quinta-feira». A frase ficou mais curta. A reserva da sala, que era o trabalho todo, desapareceu pelo caminho.

\section{Uma mensagem com um pedido claro}

Antes de escrever, decide o que o destinatário precisa de saber e que ação lhe cabe. Para a reunião fictícia, podes enviar um rascunho com um pedido concreto. «Podes confirmar se há sala na quinta-feira? Aguardo a confirmação antes de marcar.» A frase funciona mesmo quando ainda faltam a hora e o local. Quem recebe sabe o que pode resolver.

\begin{lstlisting}
Prepara um rascunho para [destinatário].
O objetivo é [decisão ou pedido].
Usa apenas os factos confirmados abaixo.
Indica o responsável e o prazo se estiverem nas notas.
Se faltar algo necessário, coloca-o como pergunta.
\end{lstlisting}

\begin{promptlabbox}{Reescrever um email que pede uma decisão}
\textbf{Objetivo e contexto:} Preparar um email fictício de até seis frases a partir da nota sobre a sala.

Escreve primeiro os factos que podes usar. Pede uma versão curta, dirigida à pessoa que pode confirmar a disponibilidade. Lê-a como se fosses essa pessoa. Consegues responder ao pedido sem adivinhar o que ficou decidido?

\textbf{Resultado esperado e validação:} O texto contém um pedido identificável, distingue a proposta da reserva e deixa explícita a informação em falta. Confirma o rascunho antes de o utilizar.

\textbf{Falhas possíveis:} Acrescentar hora, sala, responsável ou prazo ausente da nota.
\end{promptlabbox}

\section{Da reunião ao registo de decisões}

Uma agenda pode começar pelas duas decisões que precisam de ser tomadas, com o tempo e os documentos necessários. Nas notas finais, separa decisões, ações e dúvidas. Guarda o excerto ou momento de origem quando um compromisso for importante. Se a gravação ou a transcrição estiver autorizada, confere os nomes, as negações e os valores que podem alterar uma decisão.

\begin{center}
\small
\begin{tabularx}{\textwidth}{@{}p{.22\textwidth}XX@{}}
\toprule
Estado & Nota de origem (fictícia) & Registo utilizável \\
\midrule
Decisão & «A Rita confirma a disponibilidade.» & Rita verifica a sala. O prazo fica por indicar. \\
Proposta & «Talvez quinta-feira.» & Quinta-feira é uma opção. \\
Dúvida & «Ainda falta saber a lotação.» & Confirmar a lotação antes de convidar. \\
\bottomrule
\end{tabularx}
\end{center}

Se leres duas atas, usa um identificador por ação e atualiza o estado com a fonte mais recente. Uma decisão revogada precisa de uma marca explícita. Do contrário, o resumo da semana pode trazer de volta um compromisso que a equipa já cancelou. O método de contexto do capítulo~\ref{ch:contexto} ajuda a escolher os trechos a conservar.

\section{Planeamento com hipóteses à vista}

Uma lista de forças, fraquezas, oportunidades e ameaças (\emph{SWOT}) organiza uma discussão. A sua utilidade depende do que consegues confirmar. Se o rascunho disser «há procura por oficinas à noite», aponta que observação sustentaria essa frase. Pode ser o número de pedidos recebidos num período. «As pessoas gostam» deixa por explicar quem foi ouvido e quem ficou de fora.

Escolhe um objetivo e uma forma de o observar antes de pedir um plano longo. Num projeto fictício de marcações, uma opção é medir quantos pedidos ficam sem resposta após dois dias úteis. Define quando começa a contagem, que dias entram e quais os pedidos excluídos. Regista quanto tempo a equipa consegue dedicar à experiência. Se a primeira semana contrariar a hipótese, o plano precisa de um ponto onde seja possível parar ou mudar de direção.

\begin{promptlabbox}{Auditar um plano com perguntas concretas}
\textbf{Objetivo e contexto:} Examinar um plano público ou fictício com uma ação, uma alegação e uma métrica proposta.

Pede a um modelo, ou a um colega, que selecione a frase com maior impacto na decisão. Para essa frase, escreve o facto que a sustentaria, o dado disponível e a decisão que mudaria se o dado fosse outro.

\textbf{Resultado esperado e validação:} A revisão aponta um localizador, um teste ou um responsável que possa resolver a dúvida. Uma alegação por confirmar fica assinalada como tal.

\textbf{Falhas possíveis:} Transformar uma hipótese em facto, medir só o que é fácil ou usar informação interna num ambiente sem autorização.
\end{promptlabbox}

\section{De perguntas frequentes a respostas autorizadas}

As mensagens de clientes ajudam a encontrar perguntas recorrentes. A resposta deve vir da política em vigor, com as exceções e a origem à vista. Regista a data de revisão de cada resposta e o que fazer se a política não se aplicar. Na Lavandaria fictícia, o assistente pode preparar uma resposta sobre uma reclamação. Uma promessa de reembolso é uma decisão que exige autoridade e informação adicionais.

\begin{lstlisting}
Agrupa as perguntas por necessidade do cliente.
Para cada grupo, propõe uma resposta com base
na política indicada e aponta o excerto usado.
Se a política não cobrir o caso, indica a dúvida
a encaminhar e a equipa que a pode esclarecer.
\end{lstlisting}

Para avaliar uma página de ajuda, compara os rascunhos com uma amostra de perguntas que não usaste durante a escrita. Inclui casos em que a resposta precisa de ser encaminhada. Regista a proporção de perguntas cobertas, as correções necessárias e o tempo de revisão. Se automatizares a publicação ou o envio, coloca a aprovação e a verificação de permissões no ponto em que o compromisso se torna efetivo. Consulta a operação no capítulo~\ref{ch:llmops} e as decisões sobre dados no capítulo~\ref{ch:governacao}.
